\documentclass[10pt,a4paper]{amsart}
\usepackage[margin=19mm]{geometry}
\usepackage{amsmath,amssymb,mathtools}
\usepackage[scr=rsfs,cal=euler]{mathalfa}
\usepackage{xfrac}
\usepackage[unicode,hidelinks,bookmarks=false]{hyperref}
\newcommand{\ie}{i.e.\ }
\newcommand{\cf}{cf.\ }
\newcommand{\resp}{respectively\ }
\newcommand{\Subsection}{\subsection}
\providecommand{\nobreakdash}{\nobreak\hbox{-}}
\setlength{\emergencystretch}{2em}
\multlinegap0pt
\usepackage{bm}
\usepackage{bbm}
\usepackage{dsfont}
\usepackage{xspace}
\usepackage{cancel}
\usepackage{mathtools}
\usepackage{amsthm}
\makeatletter
\@ifundefined{theorem}{\newtheorem{theorem}{Theorem}}{}
\@ifundefined{lemma}{\newtheorem{lemma}[theorem]{Lemma}}{}
\makeatother
\makeatletter
\expandafter\ifx\csname CHANGED\endcsname\relax
\newenvironment{CHANGED}{\begingroup\color{blue}}{\endgroup}
\fi
\expandafter\ifx\csname MW\endcsname\relax
\newenvironment{MW}{\begingroup\color{red}}{\endgroup}
\fi
\renewcommand{\P}{{\mathbb P}}
\newcommand{\R}{{\mathbb R}}
\expandafter\ifx\csname RJK\endcsname\relax
\newenvironment{RJK}{\begingroup\color{darkgreen}}{\endgroup}
\fi
\newcommand{\expect}{\operatorname{\mathbb{E}}}
\DeclareMathOperator{\precond}{Precond}
\DeclareMathOperator{\sgn}{sgn}
\renewcommand{\vec}{\boldsymbol}
\makeatother
\title[Adaptive and non-adaptive randomized approximation of high-d]{Adaptive and non-adaptive randomized approximation of high-dimensional vectors}
\author{Robert J. Kunsch, Marcin Wnuk}
\date{Source: arXiv:2410.23067v2 (CC BY 4.0)}
\begin{document}
\maketitle

\noindent\emph{Source and licence.} Adaptive and non-adaptive randomized approximation of high-dimensional vectors. Version arXiv:2410.23067v2 (CC BY 4.0), distributed under the Creative Commons Attribution 4.0 International licence, as recorded in the arXiv licence metadata for this exact version and shown on the abstract page https://arxiv.org/abs/2410.23067v2. Full rights notice: licenses/arxiv-2410.23067-CC-BY-4.0.txt.\par\smallskip

\noindent\textit{Editorial scope.} Counted unit: the Lemma whose source label is \texttt{lem:precond} in the article (source0.tex lines 821--839, source label \texttt{lem:precond}), delivered together with the article's own complete original proof of it (source0.tex lines 840--934, one \texttt{proof} environment of 95 lines). No mathematics is written, repaired or re-derived: statement and proof are the article's own text line for line. Required context retained uncounted: none. Every definition, display and label of those lines is retained unchanged. Omitted from this package and not redistributed: 1--820, 935--1991. Nothing inside a retained excerpt is omitted or altered: every retained body is delivered exactly as the article's lines are written, so no omission receipt of any kind accompanies this package. Citations: the retained excerpts carry no citation command at all, so no bibliography entry of the article is transcribed and no reference is redistributed. Reference adaptation: none was needed and none was made; every reference inside the delivered unit resolves inside it, so no unresolved reference or citation is produced. Printed slips kept literal: no printed word, symbol, hypothesis or number of the counted statement or of its proof was altered. Layout and class: the article's own class is \texttt{article} with packages including geometry, babel, latexsym, amsfonts, amsmath, amssymb, epsfig, tabularx, amsthm, dsfont, mathrsfs, graphicx, enumerate, booktabs; the wrapper is the lane's pinned \texttt{amsart} editorial wrapper at 10pt on A4, which restates the theorem-like environments the retained text needs and adds the article's own macro definitions that the retained text needs (\texttt{\textbackslash P}, \texttt{\textbackslash R}, \texttt{\textbackslash expect}, \texttt{\textbackslash precond}, \texttt{\textbackslash sgn}, \texttt{\textbackslash vec}), with extra wrapper packages (bm, bbm, dsfont, xspace, cancel, mathtools). The delivered numbering is the unit's own. Assets: no retained span contains a figure environment, a graphics-inclusion command, a PDF-page inclusion, a table or a TikZ picture; no image file is copied into this package, the package declares no asset dependency and no image file is redistributed with it. Licence: version arXiv:2410.23067v2 of this article is distributed under the Creative Commons Attribution 4.0 International licence, as recorded in the arXiv licence metadata for this exact version and shown on the abstract page https://arxiv.org/abs/2410.23067v2. A full rights notice is delivered at licenses/arxiv-2410.23067-CC-BY-4.0.txt. Characters: every retained excerpt is pure ASCII, so the delivered engine, XeLaTeX with its default Latin Modern text font, reports no missing character.
\begin{lemma} \label{lem:precond}
    For $\vec{x} \in \R^m$ and $j^\ast \in J \subset [m]$ assume the heavy hitter condition
    \begin{equation*} \label{eq:weakHHC}
        \|\vec{x}_{J \setminus\{j^\ast\}}\|_2
            \leq \frac{|x_{j^\ast}|}{\sqrt{5}} \,.
    \end{equation*}
    Let $\gamma > 1$ and $\delta_1 \in (0,1)$.
    If we choose
    \begin{equation*}
        k := \left\lceil 36 \log\left(\frac{1+\frac{2}{5}\gamma^2}{\delta_1}\right)\right\rceil ,
    \end{equation*}
    then, with $S := \precond_{k}(\vec{x},J) \subseteq J$, we have
    \begin{equation*}
        \P\left( j^\ast \in S \quad \text{and} \quad
            \|\vec{x}_{S \setminus\{j^\ast\}}\|_2 \leq \frac{|x_{j^\ast}|}{\gamma}
            \right)
            \geq 1-\delta_1 \,.
    \end{equation*}
\end{lemma}

\begin{proof}
    The idea is that $\vec{s} \approx \sgn(x_{j^\ast}) \cdot \vec{a}_{j^\ast}$
    in the sense that likely only few entries of these two vectors differ.
    In fact, if
    \begin{equation*}
        Y_i := \sum_{j \in J \setminus\{j^\ast\}} 
                    \frac{x_j}{x_{j^\ast}a_{i,j^\ast}} \cdot a_{ij}
            > -1
    \end{equation*}
    holds then $s_i = \sgn(x_{j^\ast}) \cdot a_{i,j^\ast}$.
    Here, $Y_i$ is the sum of independent random variables with square-summable absolute bounds $\frac{|x_j|}{|x_{j^\ast}|}$. By Hoeffding's inequality we have
    \begin{align*}
        \P\bigl(s_i \neq \sgn(x_{j^\ast} a_{i,j^\ast})\bigr)
            &\leq \P(Y_i \leq -1)
            \leq \exp\left(-\frac{x_{j^\ast}^2}{2\|\vec{x}_{J\setminus\{j^\ast\}}\|_2^2}\right)
            \\
            &\leq \exp\left(-\frac{5}{2}\right) < \frac{1}{12} \,.
    \end{align*}
    It follows that
    \begin{equation*}
        \mu := \expect \left[d_H(\vec{s}, \sgn(x_{j^\ast}) \cdot \vec{a}_{j^\ast})\right]
            < \frac{k}{12} \,.
    \end{equation*}
    Since 
    $\frac{1}{2}|s_i - \sgn(x_{j^\ast}) \cdot a_{i,j^\ast}| \in \{0,1\}$
    are i.i.d.\ Bernoulli random variables,
    we may use a Chernoff bound for binomially distributed random variables,
    namely, with $\delta = \frac{k}{6\mu} - 1 \geq 1$ we find
    \begin{align*}
        \P\left(d_H(\vec{s}, \sgn(x_{j^\ast}) \cdot \vec{a}_{j^\ast})
                > \frac{k}{6}\right)
            &= \P\bigl(d_H(\vec{s}, \sgn(x_{j^\ast}) \cdot \vec{a}_{j^\ast})
                > (1+\delta)\mu\bigr) \\
            &\leq \exp\left(- \frac{\min\{\delta,\delta^2\} \cdot \mu}{3}\right)
            = \exp\left(- \frac{k/6 - \mu}{3}\right) \\
            &\leq \exp\left(-\frac{k}{36}\right) =: \alpha \,.
    \end{align*}
    This gives a guarantee for correctly identifying a subset of~$J$ that still contains the important coordinate $j^\ast$:
    \begin{align}
        \P\bigl(j^\ast \in \precond_k(\vec{x},J)\bigr)
            &\geq \P\left(d_H(\vec{s}, \sgn(x_{j^\ast}) \cdot \vec{a}_{j^\ast})
                \leq \frac{k}{6}\right)
                \nonumber \\
            &\geq 1 - \exp\left(-\frac{k}{36}\right)
            = 1 - \alpha \,.
            \label{eq:precondtrue}
    \end{align}

    Now suppose that $j^\ast \in S = \precond_k(\vec{x},J)$. Then, by the triangle inequality we further know
    \begin{equation*}
        S \subseteq \left\{ j \in J \colon \min\{d_H(\vec{a}_j,\vec{a}_{j^\ast}),d_H(\vec{a}_j,-\vec{a}_{j^\ast})\}
                \leq \frac{k}{3} \right\} =: \overline{S} \,.
    \end{equation*}
    For $j \in J \setminus\{j^\ast\}$, the quantity $d_H(\vec{a}_j,\vec{a}_{j^\ast})$ follows a symmetrical binomial distribution with expectation $\mu := \expect\left[d_H(\vec{a}_j,\vec{a}_{j^\ast})\right] = \frac{k}{2}$.
    Another Chernoff bound,
    with $\delta = \frac{1}{3}$,
    gives
    \begin{align*}
        \P\left(d_H(\vec{a}_j,\vec{a}_{j^\ast}) \leq \frac{k}{3}\right)
            &= \P\bigl(d_H(\vec{a}_j,\vec{a}_{j^\ast}) \leq (1-\delta) \mu\bigr) \\
            &\leq \exp\left(-\frac{\delta^2 \mu}{2}\right)
            = \exp\left(-\frac{k}{36}\right) .
    \end{align*}
    We can use this to estimate the mean squared norm of these coordinates:
    \begin{align*}
        \expect \|\vec{x}_{\overline{S} \setminus\{j^\ast\}}\|_2^2
            &\leq \sum_{j \in J \setminus\{j^\ast\}} \P\left(d_H(\vec{a}_j,\vec{a}_{j^\ast}) \leq \frac{k}{3} \text{ or } d_H(\vec{a}_j,-\vec{a}_{j^\ast}) \leq \frac{k}{3}\right) \cdot |x_j|^2 \\
            &\leq 2\exp\left(-\frac{k}{36}\right)
                \cdot \|\vec{x}_{J \setminus \{j^\ast\}}\|_2^2 \,.
    \end{align*}
    By Markov's inequality we find for $\gamma > 1$:
    \begin{align}
        \P\left( \|\vec{x}_{\overline{S} \setminus\{j^\ast\}}\|_2 > \frac{|x_{j^\ast}|}{\gamma}\right)
            &\leq \P\left( \|\vec{x}_{\overline{S} \setminus\{j^\ast\}}\|_2 
                    > \frac{\sqrt{5} \cdot \|\vec{x}_{J \setminus\{j^\ast\}}\|_2}{\gamma}\right)
            \nonumber \\
            &\leq \frac{2}{5} \, \gamma^2 \cdot \exp\left(-\frac{k}{36}\right) =: \beta\,.
            \label{eq:precondnorm}
    \end{align}

    Finally, combining the findings~\eqref{eq:precondtrue} and~\eqref{eq:precondnorm}, we obtain
    \begin{align*}
        \P&\left( j^\ast \in S \quad \text{and} \quad
            \|\vec{x}_{S \setminus\{j^\ast\}}\|_2 \leq \frac{|x_{j^\ast}|}{\gamma}
            \right) \\
            &\geq \P\left( j^\ast \in S \quad \text{and} \quad
            \|\vec{x}_{\overline{S} \setminus\{j^\ast\}}\|_2 \leq \frac{|x_{j^\ast}|}{\gamma}
            \right) \\
            &\geq 1 - \P\left( j^\ast \notin S\right)
                    - \P\left(\|\vec{x}_{\overline{S} \setminus\{j^\ast\}}\|_2 > \frac{|x_{j^\ast}|}{\gamma}\right) \\
            &\geq 1 - \alpha - \beta \,.
    \end{align*}
    The choice of $k$ in the lemma ensures
    $\alpha + \beta = \left(1 + \frac{2}{5} \, \gamma^2\right) \exp\left(-\frac{k}{36}\right) \leq \delta_1$.
\end{proof}

\end{document}
